\documentclass[11pt]{article}
\usepackage{amsmath,amssymb,amsthm}
\usepackage[margin=1.1in]{geometry}
\usepackage{hyperref}

\newtheorem{theorem}{Theorem}[section]
\newtheorem{lemma}[theorem]{Lemma}
\newtheorem{corollary}[theorem]{Corollary}
\newtheorem{proposition}[theorem]{Proposition}
\newtheorem{remark}[theorem]{Remark}

\title{The Galois-Closure Descent Family:\\ the non-Galois CM shapes of $R2'$ (GST Hodge Campaign, Round 37)}
\author{GST campaign continuation session (operator order: the standing finish order,\\
\texttt{research/gpt6-native-point-primitive-span-20261017}, commit after Round 36 $2c94b3c6$)}
\date{2026-11-05 (session clock)}

\begin{document}
\maketitle

\section*{Scope and the standing order}

This round strikes the face $(a')$ named by Round 36's Theorem S17(iii): the
\emph{non-Galois CM shapes} --- every CM closure pair $(\widetilde G, H)$ beyond
the abelian-cyclotomic worlds, at the committed census (degrees $4$ and $6$
complete, degree $8$ the regular and automorphism-split order-16 families).
The operator's standing order is the finish order against Question 2; per the
closing discipline (\S40) the Clay statement is \textbf{not} claimed here, and
the polarity split stands. Everything below is decided by the certificate
\texttt{galois\_closure\_descent\_round37\_checks.py} (1950 checks, 0 failures,
exit 0, deterministic double-run byte-identical, md5
\texttt{50d0b69880eb96a2bff95afe7da92c52}, $\approx$134\,s, exact arithmetic,
zero floats; the post-repair state after AUDIT-37 --- the pre-repair
certificate stood at 1892 checks, md5
\texttt{e16164e5ae0632572db5f01c93025e29}).
\emph{Provenance}: the branch identifier and the Round-36 push hash in the
byline are the receipts recorded in the campaign worklog (CX-16) and the
Round-32 record; the round artifacts live under
\texttt{hodge-campaign/round37/} on the research branch.

\section{The closure letter model}

\subsection{Setup}
Let $K$ be a CM field of degree $2g$ with Galois closure $\widetilde K$,
$G := \operatorname{Gal}(\widetilde K/\mathbb Q)$ (order $n$),
$H := \operatorname{Gal}(\widetilde K/K)$, and complex conjugation
$\rho \in Z(G)$ acting freely on the letter set
$X := G/H$ (the $2g$ embeddings $K \hookrightarrow \widetilde K$; freeness of
$\rho$ is the CM condition $\rho \notin \bigcup_g gHg^{-1}$). A \emph{CM type}
$\Phi \subseteq X$ picks one letter from each $\rho$-pair. Write
$A(K,\Phi)$ for the abelian variety of dimension $g$ with complex
multiplication by $\mathcal O_K$ and type $\Phi$, and
$\mathcal H^{p,p}_\Phi$ for the span of the letter words
$e_U = e_{u_1}\wedge\dots\wedge e_{u_{2p}}$ ($U \subseteq X$, $|U| = 2p$,
$|U \cap \Phi| = p$) in $\bigwedge^{2p}(\widetilde K \otimes H^1)$.

\begin{lemma}[Identification]\label{lem:ident}
$H^{2p}(A(K,\Phi),\mathbb Q)\otimes_{\mathbb Q}\widetilde K
\cong \bigoplus_U \widetilde K\, e_U$ compatibly with the $G$-action
$g\cdot e_\sigma = e_{g\circ\sigma}$ on the letters and the Galois action on
coefficients; consequently
\[
\operatorname{Hdg}^{2p}\!\big(A(K,\Phi)\big)
\;=\;
\Big\{\,c=(c_U)\in \widetilde K^{\,\mathcal B_\Phi}\;:\;
c_{gU}=g(c_U)\,s(g,U)\ \ \forall g\in G\,\Big\},
\]
where $\mathcal B_\Phi$ is the set of \emph{$\Phi$-balanced} words
($|U\cap\Phi|=p$) and $s(g,U)\in\{\pm1\}$ is the exterior sign of the
permutation sorting $g(U)$.
\end{lemma}

\begin{proof}
$H^1(A)\otimes\widetilde K$ is the direct sum of the $K$-eigenlines
$\widetilde K\,e_\sigma$, $\sigma\in X$; the eigenline for $\sigma$ is mapped
by $1\otimes g$ into the eigenline for $g\circ\sigma$. Choosing eigenvectors
compatibly on each $G$-orbit of letters fixes the identification up to a
scalar cocycle on the stabilizers; the cocycle multiplies the twisted action
below by a unit but never enters a trace (Lemma~\ref{lem:trace}), so the
dimensions and the product ranks computed from any compatible choice agree.
A rational class $\zeta$ has $\zeta\otimes 1=\sum_U c_U e_U$; rationality is
$(1\otimes g)(\zeta\otimes 1)=\zeta\otimes 1$, i.e.\ the displayed equivariance;
the Hodge condition $(p,p)$ is support on $\mathcal B_\Phi$.
\end{proof}

\begin{lemma}[Mixed-orbit kill]\label{lem:kill}
If $U$ is $\Phi$-balanced and $gU$ is not, then every system of
Lemma~\ref{lem:ident} has $c_U=0$. Hence the support of any rational Hodge
class lies in the set of \emph{pure-balanced} words
\[
\mathcal P_\Phi(p)=\{U : |U|=2p,\ |U\cap\Phi|=p,\ \text{and } |gU\cap\Phi|=p
\ \forall g\in G\}.
\]
\end{lemma}

\begin{proof}
The equivariance at $g$ reads $c_{gU}=g(c_U)s(g,U)$; the left side vanishes
on components outside $\mathcal B_\Phi$ by the Hodge condition, so
$g(c_U)=0$ and $c_U=0$.
\end{proof}

\subsection{The full-descent dimension law}

\begin{lemma}[Trace law]\label{lem:trace}
For $h\in G$, $h\ne e$, the trace of $h$ on $\widetilde K$ (as
$\mathbb Q$-vector space) is $0$; consequently for any subgroup
$S\le G$ and any character twist $\chi:S\to\{\pm1\}$,
\[
\dim_{\mathbb Q}\{x\in\widetilde K : h(x)=\chi(h)\,x\ \forall h\in S\}
= \frac{n}{|S|},
\]
independently of $\chi$.
\end{lemma}

\begin{proof}
By the normal basis theorem $\widetilde K\cong \mathbb Q[G]$ as a
$G$-module, and left multiplication by $h\ne e$ is a fixed-point-free
permutation of the basis. The fixed space of the twisted action is the
$\chi$-isotypic part of $\operatorname{Res}^G_S\mathbb Q[G]$, of dimension
$\frac1{|S|}\sum_{h\in S}\chi(h)\operatorname{tr}(h) = \frac{n}{|S|}$
since $\operatorname{tr}(h)=0$ for $h\ne e$ and $\chi(e)=1$.
\end{proof}

\begin{theorem}[S18, Galois-closure descent]\label{thm:s18}
Let $O$ be a $G$-orbit of pure-balanced words with stabilizer
$S=S(U)$. The systems of Lemma~\ref{lem:ident} supported on $O$ form a
$\mathbb Q$-space of dimension $|O|$; in particular
\[
\dim_{\mathbb Q}\operatorname{Hdg}^{2p}\!\big(A(K,\Phi)\big)
\;=\;\#\mathcal P_\Phi(p).
\]
\end{theorem}

\begin{proof}
A system on $O$ is determined by $c_{U_0}\in\widetilde K$ through
$c_{gU_0}=g(c_{U_0})\,s(g,U_0)$; well-definedness under $gU_0=g'U_0$
is exactly the twisted-fixed condition for $S$ on $c_{U_0}$ with the
character $s(\cdot,U_0)|_S$ (a homomorphism since $S$ stabilizes $U_0$).
Lemma~\ref{lem:trace} gives $\dim = n/|S| = |O|$. Summing over orbits
gives the dimension law. The certificate verifies this closed form against
an exact blockwise RREF solution of the equivariance system at every
committed world, type orbit, and degree (bounded
$\#\mathcal P\cdot n\le 2600$): the two algorithms agree everywhere.
\end{proof}

\subsection{The census}\label{sec:census}

\emph{Degree 4 (complete).} The 2-generation sweep of $S_4$ (all cyclic
representatives paired with all elements, conjugacy-canonicalized) finds
exactly $5$ transitive classes, matching the classical count. The
central fixed-point-free involution filter selects $\{V_4, C_4, D_4\}$:
the abelian worlds (Round 36's $m=8,12$ and $m=5$ shapes, now re-derived
through the new engine) and the \emph{non-Galois quartic} world $D_4$
($|H|=2$). $A_4$ and $S_4$ have trivial centers and are excluded.

\emph{Degree 6 (complete).} The same sweep of $S_6$ finds exactly $16$
transitive classes, matching the classical count: since every class found
is 2-generated, a hypothetical 17th transitive subgroup would contradict
the classical census, so the sweep is complete. The filter selects four
$\rho$-worlds: the abelian $C_6$ (Round 36's $m=7,9$ shapes) and the
non-abelian classes of orders $12$ ($D_6$, the generic sextic CM closure),
$24$, and $48$.

\emph{Degree 8 (committed family).} The five regular worlds (all abstract
groups of order $8$: $C_8$, $C_4\times C_2$, $C_2^3$, $D_4$, $Q_8$,
pairwise non-isomorphic by involution counts $1/3/7/5/1$ with $C_8$ vs
$Q_8$ split by exponent) and the order-16 worlds
$\langle N_{\mathrm{reg}},\tilde\alpha\rangle$ for automorphisms
$\alpha$ of each $N$ (complete in the class of transitive order-16
subgroups \emph{containing} a regular subgroup: any such $G$ has
$G\cap N_{\mathrm{reg}}=N$ and $G/N\cong C_2$, so $G=\langle
N,\tilde\alpha\rangle$ inside the holomorph). Worlds with multiple central
fixed-point-free involutions contribute one world per $\rho$. Total:
$110$ $(G,\rho)$-worlds.

\subsection{The dimension laws at the census}

\begin{proposition}[Type classification and the classical laws]\label{prop:laws}
\begin{enumerate}
\item $\Phi$ is induced from the CM subfield $k=\widetilde K^{H'}$
($H<H'<G$, $\rho\notin H'$) iff $\Phi$ is a union of right $H'$-orbits of
letters (the fibers of $G/H\to G/H'$).
\item At every \emph{primitive} type of every committed world (116 type
orbits across the census),
$\dim\operatorname{Hdg}^{2p}= \binom{g}{p}$ for all $p$; in particular
$\dim\operatorname{NS}\otimes\mathbb Q = g$ (Rosati/Weil).
\item At every type induced from a \emph{quadratic} CM subfield,
$\dim\operatorname{Hdg}^{2p}=\binom gp^2$ (the split law; the
$E^g$-worlds). The certificate checks this at every
\emph{quadratic-inducible} orbit --- any index-2 stabilizing $H'$, not only
the maximal-subfield-2 ones: 128 orbits, including the order-refined octic
split types whose maximal inducing subfield is quartic.
\item The palindrome law $\dim(2p)=\dim(2g-2p)$ for $1\le p\le g-1$ and
the top degree $\dim(2g)=1$ hold at every type.
\item Every type on a $V_4$-Galois quartic is induced (for
$\Phi=\{x,y\}$, the subgroup $\langle x^{-1}y\rangle$ stabilizes $\Phi$
since $(x^{-1}y)^2=e$ in $V_4$ and $y\ne\rho x$); likewise every type on a
$D_4$-\emph{regular} octic is induced from one of the four quartic CM
subfields. Hence $V_4$-Galois quartics and $D_4$-Galois octics admit no
simple CM abelian varieties; the computed values
$\dim\operatorname{NS}\in\{4\}$ (quartic) and $\{8\}$ (octic) are the
product-world values $B\times B'$.
\item The $Q_8$-regular octic has only primitive types (its order-4
subgroups all contain $\rho$, so it has no proper CM subfield), with the
binomial row $[4,6,4,1]$.
\end{enumerate}
\end{proposition}

Round 36's receipts are reproduced through the new engine at the abelian
worlds: the quartic span $1$, the sextic $3/9$, and the octic
$6/18/36$ rows (primitive/intermediate/split).

\section{The pairing-generation theorem}

\begin{theorem}[S19, non-Galois pairing generation]\label{thm:s19}
At every committed world $(G,\rho)$ of the census
(\S\ref{sec:census}), every CM type orbit, and every degree $2p\le 2g$:
every rational Hodge class on $A(K,\Phi)$ is a $\mathbb Q$-combination of
\emph{products of divisor classes}. Equivalently
\[
\operatorname{rank}\operatorname{span}\{d_{i_1}\wedge\dots\wedge d_{i_p}\}
=\dim_{\mathbb Q}\operatorname{Hdg}^{2p}\!\big(A(K,\Phi)\big),
\qquad\textsc{surplus}=0,
\]
the products taken with repetition (the diagonal $d\wedge d\ne0$ for
non-decomposable $d$; the iterate-basis sampler
$S_p=\operatorname{span}(S_{p-1}\wedge \mathrm{NS})$ includes them by
construction). The certificate computes \textbf{1386} (world, type,
degree) surplus entries: \textbf{0 are nonzero}.
\end{theorem}

\begin{proof}
The cup-product algebra of the complex torus $A(\mathbb C)$ is the exterior
algebra of $H^1$, so products of degree-2 classes are exterior products of
the corresponding systems; products of equivariant systems are equivariant
(cocycle multiplicativity), and products of balanced words are balanced, so
the products land in the Hodge space. The span rank is computed exactly by
integer rank maintenance over the pure-word coordinates; the certificate
additionally verifies (i) the product-equivariance relation on explicit
products, (ii) the random-class nondegeneracy
$v^{\wedge p}\ne0$ at every world (deterministic xorshift64* draws, with
bounded retries --- a single draw can be degenerate), and (iii) the
calibration battery of Proposition~\ref{prop:laws}.
\end{proof}

\begin{corollary}[S19, rational Hodge conjecture at the census]\label{cor:s19}
The rational Hodge conjecture holds, at every degree, on every CM abelian
variety whose closure pair $(\widetilde G,H)$ lies in the committed census
--- simple (the primitive type orbits, including all non-Galois primitive
shapes at degrees 4 and 6 and the committed degree-8 families), imprimitive
(the split and intermediate product worlds), non-abelian-Galois (the
$D_4$-, $Q_8$-regular and automorphism-split order-16 octics), and on
$E\times A$ over them (63 joint frames over 17 base worlds: the complete
quartic and sextic census worlds, the five regular octic worlds, and three
distinct automorphism-split order-16 bases $C_8$, $C_4\times C_2$,
$C_2^3$; per frame the first two type-orbit representatives --- the disjoint
frames calibrated by the tensor law $\dim=1+\mathrm{NS}(A)$, the coupled
frames by the fiber rule --- the full-fiber types carry
$\dim\operatorname{Hom}(E,A)=2|\Phi_A|$ and the $E^{g+1}$-law, the
straddling types carry $\dim\operatorname{Hom}=0$: no elliptic subvarieties
in the simple worlds, distinct CM fields in the split ones).
\end{corollary}

\section{The re-drawn boundary (S20)}

\begin{theorem}[S20]
\begin{enumerate}
\item[(i)] The elliptic face of $R2'$ is closed (Round 35); the abelian
non-elliptic face at the abelian-cyclotomic shapes is closed (Round 36);
\textbf{the non-Galois / non-abelian-Galois face is now closed at the
committed census} --- the merged-sector surplus at these shapes is
\emph{zero} at every degree, type, and arity covered.
\item[(ii)] The residual of $(a')$ beyond the census: degree-8 closure
groups of order $\ge24$ or without regular subgroups, and degrees $\ge10$
--- the family-transport law (\S45) carries the statement shape; no claim.
\item[(iii)] The remaining face of $R2'$: $(b)$ the non-split Weil
fourfold world ($\mathrm{NS}(A)_{\mathbb Q}=\mathbb Q\theta$, the
2-dimensional Weil plane, the committed obstruction receipts of Round 34).
\item[(iv)] $R3$ (the $K$-condition at general input beyond the channels)
and $R4$ (the transcendental fold-birth with the $(p,p)$-sensitive gate)
stand untouched.
\end{enumerate}
\end{theorem}

\section*{The certificate}

\texttt{galois\_closure\_descent\_round37\_checks.py}: V1 the census
(completeness receipts 5/5 and 16/16 against the classical counts; the
degree-8 family as in \S\ref{sec:census}); V2 the model at every world
(exterior cocycle law verified on all $(g,h,U)$ at degrees 2 and 4; the
mixed-orbit kill live at 224 type-orbits; the two-algorithm dimension
agreement at all 110 worlds); V3 the crux (the classifier, the calibration
battery of Proposition~\ref{prop:laws}, \textsc{surplus}$=0$ at 1386
entries, the random-class and product-equivariance receipts); V4 the
$E\times A$ frames; V5 the operator's calibration anchor (arity totals
$4/8/24/64/208$ by two algorithms at every $m\le5$, the $m=6$ control $688$
by both, $\det(I+C_L)=1-(-1)^L$ by exact Bareiss for $L=2..12$); V6 the
hostile battery (basis-deletion, wrong-sign cocycle flip, the type-validity
battery --- double-letter, missing-pair, duplicate-letter, empty, at a
canary world and adversarially at every census world --- the
unbalanced-word injection into the filter's input list, garbage liveness,
the census pair-generation and generator-deletion --- every
control fires with the predicted failure) plus the mechanical vacuity and
self-comparison scans. Two defects were caught and repaired during
authoring by the certificate's own calibrations (the coupled-frame
closure late-binding and the single-draw degeneracy); both are documented
in the file. AUDIT-37 then found and repaired thirteen further findings in
one batched round (the S3b one-directional-canary gap closed by the
adversarial type battery; the split-law domain widened to all 128
quadratic-inducible orbits; the E$\times$A frame census extended to 17
base worlds; the degree-2 spans made genuine; the H4/H6 controls made
load-bearing; the Lean PART-A label corrected); the closing re-check
stands at 1950/0/exit 0, deterministic.

\section*{Boundary on every artifact}

Q1 (the framework conjecture on the committed world model) remains proven
and certified, unchanged. Q2 (the Clay statement) is \textbf{not claimed};
the closing string is \textbf{not emitted}; the polarity split stands. The
falsifiable line after this round: $(b)$ the non-split Weil fourfold;
$R3$; $R4$; the transport family of $(a')$ beyond the census.

\end{document}
